% Fusión nueva: conserva original y prototipo; usa el núcleo real de póster.
% Tres tiempos visuales: introducción, producto temático con síntesis, conclusión.
% La curva es un cálculo determinista hipotético, no un registro de personas.
\documentclass[spanish,final,xcolor={dvipsnames,table}]{beamer}
\def\itescapostertitle{Muestreo para decidir}
% Adaptador local de póster ITESCA. Requiere clase beamer; no es un article.
% Hereda el núcleo real sin copiarlo ni modificar sus configuraciones globales.
% Hook por basename de LaTeX: después de config_poster y antes de imports/cmd.
\providecommand{\itescapostertitle}{Estadística aplicada para la administración}
\providecommand{\itescaposterauthor}{Martín Jonathan de la Cruz Muñoz}
\providecommand{\itescaposterstudentid}{26130503}
\providecommand{\itescaposterdate}{25 de octubre de 2026}
\providecommand{\itescapostercourse}{Estadística aplicada para la administración}
\providecommand{\itescaposterteacher}{Dra. Yoana Elizeth Tauitmes Delgado}
\providecommand{\itescaposterwidth}{29.7}
\providecommand{\itescaposterheight}{42}
\providecommand{\itescapostercolumns}{1}
\providecommand{\itescaposterfontsize}{12}
\def\documenttitle{\itescapostertitle}
\AddToHookNext{file/config_poster.tex/after}{%
  \def\posterwidth{\itescaposterwidth}
  \def\posterheight{\itescaposterheight}
  \def\colnumbers{\itescapostercolumns}
  \def\colseparation{0.025\paperwidth}
  \def\documentfontsize{\itescaposterfontsize}
  \def\fontdocument{lmodern}
  \def\stylecitereferences{custom}
  \definecolor{itescaPosterRed}{HTML}{8F1117}
  \definecolor{itescaPosterAccent}{HTML}{D70F18}
  \definecolor{itescaPosterInk}{HTML}{181818}
  \definecolor{itescaPosterPaper}{HTML}{F4F7F7}
  \definecolor{itescaPosterTeal}{HTML}{006B67}
  \setbeamercolor{normal text}{fg=itescaPosterInk,bg=white}
  \setbeamercolor{structure}{fg=itescaPosterRed}
  \setbeamercolor{headline}{fg=itescaPosterInk,bg=white}
  \setbeamercolor{headline rule}{bg=itescaPosterAccent}
  \setbeamercolor{block title}{fg=itescaPosterRed,bg=white}
  \setbeamercolor{block body}{fg=itescaPosterInk,bg=itescaPosterPaper}
  \setbeamerfont{headline title}{size=\fontsize{24}{27},series=\bfseries}
  \setbeamerfont{headline author}{size=\fontsize{10.5}{12}}
  \setbeamerfont{headline institute}{size=\fontsize{10}{12}}
  \setbeamerfont{block title}{size=\fontsize{14}{16},series=\bfseries}
  \setbeamerfont{block body}{size=\fontsize{12}{14.2}}
  \setbeamerfont{footline}{size=\fontsize{8.5}{10}}
}
\AddToHookNext{file/imports_poster.tex/after}{%
  \usepackage{fontspec}
  \setmainfont{Arial}\setsansfont{Arial}
  \usefonttheme{professionalfonts}
  \usepackage[authoryear,round]{natbib}
  \setcitestyle{citesep={;}}
  \usepackage{tikz}
  \usetikzlibrary{arrows.meta,positioning,calc,babel}
}
\makeatletter
\def\input@path{{aulaTEX/base/Plantilla-Informe/}}
\makeatother
% Template:     Poster LaTeX
% Documento:    Núcleo del template
% Versión:      DEV
% Codificación: UTF-8
%
% Autor: Pablo Pizarro R.
%        pablo@ppizarror.com
%
% Manual template: [https://latex.ppizarror.com/poster]
% Licencia MIT:    [https://opensource.org/licenses/MIT]

% -----------------------------------------------------------------------------
% CONFIGURACIONES
% -----------------------------------------------------------------------------
% Definiciones previas
\input{src/defs}

% Ajustes usuario
\input{src/config}
\input{src/config_poster}

% -----------------------------------------------------------------------------
% IMPORTACIÓN DE LIBRERÍAS
% -----------------------------------------------------------------------------
\input{src/env/imports_poster}

% -----------------------------------------------------------------------------
% IMPORTACIÓN DE FUNCIONES Y ENTORNOS
% -----------------------------------------------------------------------------
\input{src/cmd/core}
\input{src/cmd/math}
\input{src/cmd/equation}
\input{src/cmd/image}
\input{src/cmd/title}
\input{src/cmd/other}
\input{src/cmd/column}
\input{src/cmd/poster}
\input{src/env/environments}

% -----------------------------------------------------------------------------
% IMPORTACIÓN DE ESTILOS
% -----------------------------------------------------------------------------
\input{src/style/code}
\input{src/style/other}

% -----------------------------------------------------------------------------
% CONFIGURACIÓN INICIAL DEL DOCUMENTO
% -----------------------------------------------------------------------------
\input{src/cfg/init_poster}

% -----------------------------------------------------------------------------
% CONFIGURACIÓN DE PÁGINA Y ENCABEZADOS
% -----------------------------------------------------------------------------
\input{src/cfg/page_poster}

% -----------------------------------------------------------------------------
% CONFIGURACIONES FINALES
% -----------------------------------------------------------------------------
\input{src/cfg/final_poster}

% Marco local de ancho útil explícito; evita la fila exterior de columnas vacías
% del entorno poster. Usa el mismo núcleo beamerposter, cabecera, pie y bloques.
\setbeamersize{text margin left=\sepwidth,text margin right=\sepwidth}
\newenvironment{itescaposter}[1][]{%
  \begin{frame}[fragile,t,environment=itescaposter,#1]%
}{\end{frame}}
\author{\texorpdfstring{\itescaposterauthor\quad Matrícula: \itescaposterstudentid}{\itescaposterauthor; Matrícula: \itescaposterstudentid}}
\institute{Maestría en Gestión Administrativa\\\itescapostercourse\\\itescaposterteacher}
\logoleft{width=3.4cm}{ITESCA/assets-itesca/web/logo-itesca-contorno.png}
\footercontent{\textbf{VARIANTE LOCAL DE MUESTRA}\quad GGEA01\quad\itescaposterdate}
\hypersetup{hidelinks,breaklinks=true,pdftitle={\itescapostertitle},pdfauthor={\itescaposterauthor},pdfsubject={Infografía alternativa local de muestreo; sin envío a plataforma}}
% Espaciado local de bloques. El núcleo aporta beamerposter, columnas, cabecera,
% pie, comandos matemáticos, imágenes, tablas y tamaño beamerposter.
\setbeamertemplate{block begin}{%
  \begin{beamercolorbox}[wd=\linewidth,sep=0.09cm]{block title}%
    \usebeamerfont{block title}\insertblocktitle
  \end{beamercolorbox}%
  \begin{beamercolorbox}[wd=\linewidth,sep=0.18cm]{block body}%
    \usebeamerfont{block body}\justifying
    \setlength{\parindent}{0pt}\setlength{\parskip}{0.35em}%
}
\setbeamertemplate{block end}{\end{beamercolorbox}\vspace{0.22cm}}
\newenvironment{itescaposterpanel}[1]{%
  \setbeamercolor{block title}{fg=itescaPosterTeal,bg=white}%
  \setbeamerfont{block title}{size=\fontsize{13}{15},series=\bfseries}%
  \begin{block}{#1}%
}{\end{block}}

\usepackage{pgfplots}
\pgfplotsset{compat=1.18}
\setbeamerfont{headline title}{size=\fontsize{30}{33},series=\bfseries}
\setbeamerfont{block body}{size=\fontsize{12}{14.5}}
\setbeamerfont{block title}{size=\fontsize{15}{17},series=\bfseries}
\renewcommand{\insertfootercontent}{\textbf{Actividad 1: Infografía}\quad Estadística aplicada para la administración\quad 25 de octubre de 2026}
\hypersetup{pdftitle={Muestreo para decidir},pdfsubject={Actividad 1 - Estadística aplicada para la administración}}
\newlength{\fusionpanelwidth}
\setlength{\fusionpanelwidth}{0.485\colwidth}
\newcommand{\fusionimagepath}{ITESCA/maestria-en-gestion-administrativa/estadistica-aplicada-para-la-administracion-mga/assets-estadistica-aplicada-para-la-administracion/poster/ilustracion-muestreo-administrativo.png}
% Cabecera local compacta: conserva identidad y datos sin consumir el área analítica.
\setbeamertemplate{headline}{%
  \begin{beamercolorbox}[wd=\paperwidth,leftskip=0.75cm,rightskip=0.75cm,sep=0.14cm]{headline}%
    \vspace{0.20cm}
    \begin{minipage}[c]{0.15\paperwidth}
      \includegraphics[width=3.2cm]{ITESCA/assets-itesca/web/logo-itesca-contorno.png}
    \end{minipage}%
    \begin{minipage}[c]{0.77\paperwidth}
      \centering
      {\fontsize{28}{31}\selectfont\bfseries\color{itescaPosterInk} Muestreo para decidir\par}
      \vspace{0.18cm}
      {\fontsize{10.5}{12.5}\selectfont
      Martín Jonathan de la Cruz Muñoz\quad Matrícula: 26130503\par
      Maestría en Gestión Administrativa\par
      Estadística aplicada para la administración\par
      Dra. Yoana Elizeth Tauitmes Delgado\par}
    \end{minipage}\par\vspace{0.10cm}
  \end{beamercolorbox}%
  {\color{itescaPosterAccent}\hrule height 1.5pt}%
}
\setbeamertemplate{block end}{\end{beamercolorbox}\vspace{0.16cm}}
\begin{document}
\templatePagecfg\templateFinalcfg
\begin{itescaposter}
\setstretch{1}\fontsize{12}{14.5}\selectfont\setlength{\parindent}{0pt}
\begin{tikzpicture}[remember picture,overlay]
  \fill[itescaPosterRed] (current page.north west) rectangle ([xshift=0.22cm]current page.south west);
\end{tikzpicture}
\begin{columns}[T,totalwidth=\colwidth]
\begin{column}{0.545\colwidth}
\begin{block}{Introducción}
{\color{itescaPosterTeal}\bfseries Una parte de la población, una decisión con límites.}

El muestreo permite conocer un conjunto sin observarlo completo. Ahorra tiempo y costo, pero sus conclusiones dependen de la selección, la cobertura y la respuesta. Una base grande puede describir mal a quienes interesa conocer \citep{holmesBusiness2023,bradley2021}.
\end{block}
\end{column}
\begin{column}{0.435\colwidth}
\begingroup
\def\marginimagetop{0}\def\marginimagebottom{0}\def\captionlessmarginimage{0}
\insertimage{\fusionimagepath}{width=12cm,height=4cm}{}
\endgroup
\end{column}
\end{columns}

\begin{block}{Del universo al dato: elegir antes de contar}
\centering
\begin{tikzpicture}[
  box/.style={draw=itescaPosterTeal,fill=white,rounded corners=2pt,minimum height=1.35cm,text width=7.8cm,align=center,font=\fontsize{12}{14}\selectfont},
  >=Stealth,every path/.style={line width=1pt}]
\node[box] (p) {\textbf{Población}\par Clientes del periodo que interesa conocer};
\node[box,right=0.9cm of p] (f) {\textbf{Marco muestral}\par Registro que permite seleccionarlos};
\node[box,right=0.9cm of f] (s) {\textbf{Muestra}\par Clientes seleccionados para observar};
\draw[->,itescaPosterTeal] (p.east)--(f.west);
\draw[->,itescaPosterTeal] (f.east)--(s.west);
\end{tikzpicture}
\end{block}

\begin{columns}[T,totalwidth=\colwidth]
\begin{column}{\fusionpanelwidth}
\begin{itescaposterpanel}{01 / Diseños probabilísticos}
Probabilidades conocidas y positivas de inclusión; \textbf{no siempre iguales}.

\textbf{Aleatorio simple.} Todas las muestras de igual tamaño tienen la misma probabilidad: sortear clientes de un registro completo.

\textbf{Estratificado.} Seleccionar al azar dentro de todos los estratos: clientes por canal de compra.

\textbf{Sistemático.} Inicio aleatorio y cada $k$ unidades: revisar cada décima factura; evitar periodicidades.

\textbf{Conglomerados.} Sortear grupos, como sucursales, y estudiar sus unidades. Ahorra desplazamientos; la semejanza interna puede reducir precisión \citep{holmesBusiness2023}.
\end{itescaposterpanel}
\end{column}
\begin{column}{\fusionpanelwidth}
\begin{itescaposterpanel}{02 / Selección no probabilística}
No permite conocer las probabilidades de inclusión bajo un diseño probabilístico.

\textbf{Conveniencia.} Entrevistar a personas disponibles: compradores presentes en una tienda.

\textbf{Autoselección.} Responder por iniciativa propia: encuesta abierta en redes sociales.

Útil para explorar; no atribuirle el margen de una muestra aleatoria sin supuestos adicionales.

{\color{itescaPosterRed}\bfseries Más respuestas no corrigen por sí solas el sesgo.} Sortear dentro de un marco incompleto tampoco cubre automáticamente a toda la población \citep{bradley2021}.
\end{itescaposterpanel}
\end{column}
\end{columns}

\begin{columns}[T,totalwidth=\colwidth]
\begin{column}{0.545\colwidth}
\begin{itescaposterpanel}{03 / Tamaño y precisión}
\centering
\begin{tikzpicture}
\begin{axis}[
  width=14.4cm,height=6.1cm,
  xmin=50,xmax=1100,ymin=0,ymax=11,
  axis lines=left,axis line style={itescaPosterInk},
  xlabel={Respuestas completas ($n$)},
  ylabel={Margen (puntos porcentuales)},
  xlabel style={font=\fontsize{10.5}{12}\selectfont},
  ylabel style={font=\fontsize{10.5}{12}\selectfont},
  tick label style={font=\fontsize{10}{12}\selectfont},
  xtick={100,385,1000},ytick={0,2,4,6,8,10},
  grid=major,grid style={gray!18},clip=false]
\addplot[domain=80:1100,samples=140,very thick,color=itescaPosterTeal]{98/sqrt(x)};
\addplot[only marks,mark=*,mark size=2.6pt,color=itescaPosterRed] coordinates {(100,9.8) (385,4.994542475) (1000,3.099032107)};
\node[anchor=west,font=\fontsize{10.5}{12}\selectfont,text=itescaPosterRed] at (axis cs:135,9.8) {100: 9.80 pp};
\node[anchor=south west,font=\fontsize{10.5}{12}\selectfont,text=itescaPosterRed] at (axis cs:410,5.25) {385: 4.99 pp};
\node[anchor=north east,font=\fontsize{10.5}{12}\selectfont,text=itescaPosterRed] at (axis cs:1025,2.7) {1000: 3.10 pp};
\end{axis}
\end{tikzpicture}
\par\raggedright\fontsize{10.5}{12.5}\selectfont
\textbf{Escenario hipotético, no datos observados.} $e=1.96\sqrt{0.5(0.5)/n}$: 95\%, muestreo aleatorio simple y población grande respecto de la muestra. La curva muestra menor variación; \textbf{no demuestra ausencia de sesgo} \citep{holmesBusiness2023}.
\end{itescaposterpanel}
\end{column}
\begin{column}{0.435\colwidth}
\begin{itescaposterpanel}{04 / Planificar respuestas}
\begin{minipage}{0.47\linewidth}
{\color{itescaPosterRed}\fontsize{45}{48}\selectfont\bfseries 385}
\end{minipage}%
\begin{minipage}{0.50\linewidth}
\bfseries respuestas completas\par\normalfont 95\% de confianza\par margen de 5 pp
\end{minipage}
\[
n_0=\frac{1.96^2(0.5)(0.5)}{0.05^2}=384.16
\]
Redondear hacia arriba. Se usa $p=0.5$ como máxima variabilidad en esta planificación \citep{holmesBusiness2023}.

Con respuesta prevista de 80\%:
\[
\lceil385/0.80\rceil=\boxed{482}\ \text{invitaciones}.
\]
No garantiza 385 respuestas ni corrige no respuesta diferencial. Otros diseños necesitan ajustes.

\textbf{Para contrastes:} definir también potencia y efecto mínimo detectable; esta fórmula no los sustituye \citep{nistHandbook}.
\end{itescaposterpanel}
\end{column}
\end{columns}

\begin{block}{Aplicaciones: del porcentaje a una acción defendible}
\begin{columns}[T,onlytextwidth]
\begin{column}{0.485\linewidth}
\textbf{Administración / calidad de facturación.} Seleccionar facturas con inicio aleatorio, estimar errores y orientar la auditoría. Los errores detectados no demuestran sus causas.
\end{column}
\begin{column}{0.485\linewidth}
\textbf{Mercado / satisfacción por canal.} Seleccionar al azar dentro de tienda física y venta digital, y ponderar según el diseño. Escuchar solo a seguidores digitales excluye otros compradores.
\end{column}
\end{columns}
\medskip
La comparación conduce a una conclusión práctica: \textbf{definir población, marco y selección antes de ampliar la encuesta}. Así se reconoce qué representa el porcentaje y qué permanece fuera \citep{bradley2021,holmesBusiness2023}.
\end{block}

\begin{block}{Conclusiones}
Una decisión responsable revisa cobertura, selección y no respuesta antes de celebrar el tamaño de la base. Conviene invertir primero en datos pertinentes y después en cantidad: el ahorro del muestreo solo es útil si la gerencia conoce los límites de la evidencia.
\end{block}

\begin{block}{Referencias}
\fontsize{9.5}{11.1}\selectfont\setstretch{1}\raggedright
\urlstyle{same}\Urlmuskip=0mu plus 1mu\setlength{\emergencystretch}{2em}
\renewcommand{\bibsection}{}\setlength{\bibsep}{2pt}\setlength{\bibhang}{0.65cm}
\begin{thebibliography}{3}
\bibitem[Bradley et al.(2021)]{bradley2021}
Bradley, V. C., Kuriwaki, S., Isakov, M., Sejdinovic, D., Meng, X.-L., \& Flaxman, S. (2021). Unrepresentative big surveys significantly overestimated US vaccine uptake. \textit{Nature, 600}(7890), 695--700. \url{https://doi.org/10.1038/s41586-021-04198-4}
\bibitem[Holmes et al.(2023)]{holmesBusiness2023}
Holmes, A., Illowsky, B., \& Dean, S. (2023). \textit{Introductory business statistics} (2.ª ed.). OpenStax. \url{https://openstax.org/details/books/introductory-business-statistics-2e}
\bibitem[NIST y SEMATECH(2003)]{nistHandbook}
National Institute of Standards and Technology \& SEMATECH. (2003). \textit{NIST/SEMATECH e-Handbook of statistical methods}. \url{https://doi.org/10.18434/M32189}
\end{thebibliography}
\end{block}
{\fontsize{8}{9.5}\selectfont\raggedright
\textbf{Apoyo de IA y responsabilidad académica.} Se utilizó IA como apoyo para la corrección y organización del contenido. El autor declara haber realizado y validado la investigación documental, el análisis y los resultados. Imagen: ilustración conceptual generada con IA, sin datos empíricos.\par}
\end{itescaposter}
\end{document}
